\documentclass[11pt]{article}
\usepackage[T1]{fontenc}
\usepackage[margin=1in]{geometry}
\usepackage{enumitem}
\usepackage{hyperref}
\setlength{\parindent}{1.5em}
\setlength{\parskip}{6pt}
% =====================================================================
%  EDIT YOUR DETAILS HERE — this ONE file feeds every abstract & deck.
%  (Change a value, recompile; all 36 documents update.)
% =====================================================================
\newcommand{\seminarName}{Aflah Muhammed P}      % <-- your full name (guessed from your email; please verify)
\newcommand{\seminarRoll}{TVE00CS000}            % <-- your university register/roll number
\newcommand{\seminarClassRoll}{00}               % <-- your class roll no (the short one)
\newcommand{\seminarGuide}{Prof.\ <Guide Name>}  % <-- your guide
\newcommand{\seminarDept}{Department of Computer Science and Engineering}
\newcommand{\seminarCollege}{College of Engineering Trivandrum}
\newcommand{\seminarShortCollege}{CET}           % shown in the slide footer, e.g. "Aflah Muhammed P (CET)"
\newcommand{\seminarShortTitle}{B.Tech Seminar}  % middle of the slide footer
\newcommand{\seminarDate}{July 31, 2026}
% ---------------------------------------------------------------------
% Optional: put a logo image named  cet_logo.png  in this folder and it
% will appear on every title slide automatically. If absent, it is skipped.
% =====================================================================

\begin{document}
\begin{center}
{\LARGE\bfseries Search-R1: Reinforcement Learning for Search-Augmented Reasoning}\\[1.1em]
{\large\bfseries Seminar Abstract}\\[0.6em]
{\normalsize \seminarDate}
\end{center}
\vspace{0.6em}
\section*{Abstract}
Large language models encode broad knowledge in their parameters, yet they reason over a static and quickly outdated snapshot, producing factual errors on knowledge-intensive questions. Retrieval-augmented generation (RAG) supplies external evidence, but conventional pipelines retrieve at fixed points and cannot decide, mid-reasoning, when or what to look up. Agentic prompting methods bolt search onto a frozen model without any learning signal, while supervised tool-use traces are costly to annotate and generalize poorly. As questions become multi-hop, the model must interleave thinking with several targeted searches, a behavior that fixed pipelines and prompt engineering struggle to elicit reliably.

This seminar presents \emph{Search-R1}, a reinforcement-learning framework that trains an LLM to interleave step-by-step reasoning with self-issued search-engine queries and real-time retrieval across multiple turns. The policy is optimized end-to-end with a simple outcome-based reward (exact-match correctness), while a retrieved-token loss-masking scheme excludes externally retrieved text from the policy gradient to stabilize training; both PPO and GRPO are supported. Evaluated on seven single- and multi-hop QA benchmarks with Qwen2.5 models over an E5 retriever and a Wikipedia corpus, Search-R1 improves over strong RAG baselines by roughly 41\% (7B) and 20\% (3B). The talk covers the interaction format, the reward and masking design, the system architecture, the main results, and the method's current limitations and future directions.
\section*{Key Reference Papers}
\begin{enumerate}
  \item B. Jin, H. Zeng, Z. Yue, et al., ``Search-R1: Training LLMs to Reason and Leverage Search Engines with Reinforcement Learning,'' arXiv:2503.09516, 2025.
  \item DeepSeek-AI, ``DeepSeek-R1: Incentivizing Reasoning Capability in LLMs via Reinforcement Learning,'' arXiv:2501.12948, 2025.
  \item X. Li, G. Dong, J. Jin, et al., ``Search-o1: Agentic Search-Enhanced Large Reasoning Models,'' arXiv:2501.05366, 2025.
\end{enumerate}
\vspace{0.8em}
\noindent\textbf{Submitted By}\\
\seminarName\\
\seminarRoll

\vspace{0.6em}
\noindent\textbf{Guide}\\
\seminarGuide
\end{document}
